\chapter{Company Presentation}\label{ch:company}

\section{OCP Group and its field of activity}
OCP (Office Ch\'erifien des Phosphates) was founded in 1920 and is a leading global producer of phosphate-based plant nutrition products, holding about 68\% of the world's estimated phosphate reserves. Mining began at Khouribga in 1921 and extended to Youssoufia (1931) and Benguerir, while chemical production started at Safi (1965) and Jorf Lasfar (1986); a slurry pipeline linking Khouribga to Jorf Lasfar was inaugurated in 2014 \citep{wikiocp}. For the first half of 2026 the group reported revenue of MAD 48.4~billion and an EBITDA margin of 28\%, despite weaker fertilizer demand and a sharp rise in sulfur prices \citep{mwn2026}.

\begin{figure}[!htbp]\centering
\begin{tikzpicture}[node distance=5mm,
 st/.style={draw=OcpGreen,fill=LightGray,rounded corners=2pt,minimum height=12mm,text width=3.0cm,align=center,font=\small},
 lab/.style={font=\scriptsize\itshape,text=OcpBlue,align=center,text width=3.0cm},
 arr/.style={-{Latex[length=2mm]},thick,draw=OcpGreen}]
\node[st](a){\textbf{Mining}\\ phosphate rock};
\node[st,right=of a](b){\textbf{Beneficiation and transport}\\ washing, slurry pipeline};
\node[st,right=of b](c){\textbf{Chemical platforms}\\ phosphoric acid, fertilizers};
\node[st,right=of c](d){\textbf{Marketing and export}\\ ports, customers};
\node[lab,below=1mm of a]{Khouribga, Gantour (Youssoufia, Benguerir), Boucraa};
\node[lab,below=1mm of b]{Khouribga--Jorf Lasfar pipeline (2014)};
\node[lab,below=1mm of c]{Jorf Lasfar, Safi\\ \textcolor{OcpGreen}{\textbf{TSP lines}}};
\node[lab,below=1mm of d]{Five continents};
\draw[arr](a)--(b);\draw[arr](b)--(c);\draw[arr](c)--(d);
\end{tikzpicture}
\caption{Simplified OCP value chain; the project belongs to the chemical stage.}\label{fig:valuechain}
\end{figure}

Figure~\ref{fig:valuechain} places the internship in the chemical stage, where phosphoric acid and ground phosphate are converted into granular fertilizers. The data used by the project describe a TSP line: acid and phosphate feeds, reaction and passage tanks, slurry spraying, a dryer fed by fuel oil and hot air, and a recycle stream.

\section{Organisational structure}
OCP is organised along its value chain: mining sites, chemical platforms, and commercial and logistics activities, supported by transversal functions such as finance, human resources, maintenance and information systems. Figure~\ref{fig:org} gives a simplified functional view of the units that interact with this project; it is a working representation built from the group's published value chain and the roles met during the internship, not the official organisation chart.

\begin{figure}[!htbp]\centering
\begin{tikzpicture}[node distance=4mm and 4mm,
 top/.style={draw=OcpBlue,fill=OcpBlue,text=white,rounded corners=2pt,minimum height=9mm,text width=3.6cm,align=center,font=\small\bfseries},
 mid/.style={draw=OcpBlue,fill=white,rounded corners=2pt,minimum height=9mm,text width=2.75cm,align=center,font=\small},
 hl/.style={draw=OcpGreen,fill=LightGray,rounded corners=2pt,minimum height=9mm,text width=2.75cm,align=center,font=\small},
 ln/.style={thick,draw=OcpBlue}]
\node[top](g){OCP Group};
\node[mid,below left=8mm and 12mm of g](m){Mining};
\node[hl,below=8mm of g](c){Chemical platform: TSP production unit};
\node[mid,below right=8mm and 12mm of g](k){Marketing and logistics};
\node[mid,below left=8mm and 1mm of c](p){Process and production engineering};
\node[mid,below=8mm of c](a){Automation and control (PLC/DCS)};
\node[mid,below right=8mm and 1mm of c](i){Industrial IT/OT and data};
\draw[ln](g)--(m);\draw[ln](g)--(c);\draw[ln](g)--(k);
\draw[ln](c)--(p);\draw[ln](c)--(a);\draw[ln](c)--(i);
\end{tikzpicture}
\caption{Simplified functional structure around the host production unit.}\label{fig:org}
\end{figure}

\section{Professions involved in the project}
The prototype sits at the intersection of several professions (Table~\ref{tab:roles}), each of which is an interface of the design: process engineers define what the signals represent, automation engineers own the control layer that the prototype never touches, industrial IT/OT teams govern the network path, and data and software engineers maintain the models and the service.

\begin{table}[!htbp]\centering\small\zebra
\caption{Professions met in the project and their link with the prototype.}\label{tab:roles}
\begin{tabularx}{\textwidth}{L{0.27\textwidth}YY}\toprule
\thead\hd{Profession} & \hd{Responsibility} & \hd{Contribution expected by the project}\\\midrule
Process and production & Operating conditions, product quality & Signal semantics, acceptable error, operating-event records\\
Automation and control & PLC/DCS logic, interlocks, tags & Tag list and Node IDs; read-only integration with no write path\\
Industrial IT and OT & Network zones, OPC UA endpoint, certificates & Gateway access, security policy, firewall route, host placement\\
Data and software engineering & Features, models, containers, storage, dashboard & Feature contract, inference service, tests, monitoring interface\\
Maintenance & Sensors and communications & Sensor health and communication status\\\bottomrule
\end{tabularx}
\end{table}
